\chapter{Theoretical Foundations}\label{chap:fondements}

\section{Introduction}
\TODO{Chapter introduction}

\section{Data Mining Process}
% KDD / CRISP-DM process: selection, preprocessing, transformation, data mining, interpretation.
% Show where each stage of the project fits (a draw.io diagram is possible).
\TODO{Process overview and its connection to the following chapters}

\section{Supervised Classification and Selected Models}
% For EACH model: principle, key hyperparameters, strengths and limitations, and why it was selected for this project.
% Use LaTeX formulas; include only those relevant to the experiments.
\subsection{Baseline Classifier (Zero-R)}
\TODO{Principle and role as a lower-bound baseline}

\subsection{K-Nearest Neighbors}
\TODO{Distance, choice of K, standardization, weighting}

\subsection{Decision Tree and Random Forest}
\TODO{Gini index, entropy, pruning, bagging, feature importance}

\subsection{Naive Bayes}
\TODO{Bayes' theorem, independence assumption, smoothing}

\subsection{Support Vector Machines}
\TODO{Margin, kernels, parameters C and gamma}

\section{Model Evaluation}
% Confusion matrix, precision, recall, F1-score, macro-F1, balanced accuracy, one-vs-rest ROC-AUC.
% Explain why accuracy alone can be misleading when one class is extremely rare.
\TODO{Definitions with formulas (\texttt{equation}) and justified metric selection}

\section{Cross-Validation and Hyperparameter Tuning}
% Stratified K-Fold, grouped cross-validation, GridSearchCV / RandomizedSearchCV (K >= 5),
% separation between model tuning (cross-validation) and final testing.
\TODO{Principles and justification of the evaluation protocol}

\section{Statistical Tests}
% Kruskal-Wallis (compare Bytes across the four actions), chi-square test of independence and Cramér's V
% (zero or non-zero NAT Source Port against Action), Spearman correlation.
% For each test: hypotheses, test statistic, conditions of use, and why it is appropriate here.
\TODO{Three tests: definition, explanation, and role in the project}

\section{Class Imbalance, Duplicates, and Data Leakage}
% Rare classes, exact duplicates, conflicting labels, data leakage (features known only after the decision),
% grouped splitting to prevent identical vectors from appearing in both training and test sets.
\TODO{Concepts and risks, with references}

\section{Conclusion}
\TODO{Chapter conclusion and transition to Chapter~\ref{chap:donnees}}